\subsection{DEFINITION AND GENERAL CRITERIA}

DEFINITION 1. Let X be a topological space and Y a uniform space. A subset H of F(X; Y) is said to be equicontinuous at a point \(x_{0} \in X\) if, for each entourage V of Y, there is a neighbourhood U of \(x_{0}\) in X such that \((f(x_{0}),f(x))\in\mathbf{V}\) for all \(x\in\mathbf{U}\) and all \(f\in\mathbf{H}\) . H is said to be equicontinuous if it is equicontinuous at every point of X.

DEFINITION 2. Let X and Y be two uniform spaces. A subset H of \(\mathcal{F}(X;Y)\) is said to be uniformly equicontinuous if, for each entourage V of Y, there exists an entourage U of X such that we have \((f(x), f(x')) \in V\) whenever \((x, x') \in U\) and \(f \in H\) .

A family \((f_{t})_{t\in I}\) of mappings of X into Y is said to be equicontinuous at a point \(x_0\) (resp. equicontinuous, uniformly equicontinuous) if the set of the \(f_{t}\) is equicontinuous at \(x_0\) (resp. equicontinuous, uniformly equicontinuous).

It is clear that if \(\mathbf{H}\subset\mathcal{F}(\mathbf{X};\mathbf{Y})\) is equicontinuous at \(x_{0}\) , then each \(f\in H\) is continuous at \(x_{0}\) ; if H is equicontinuous, then each \(f\in H\) is continuous on X, i.e.~\(\mathbf{H}\subset\mathcal{C}(\mathbf{X};\mathbf{Y})\) . Likewise, if H is uniformly equicontinuous (X being a uniform space), every \(f\in H\) is uniformly continuous on X. It is also clear that a uniformly equicontinuous set of mappings is equicontinuous; but a set of uniformly continuous mappings can be equicontinuous without being uniformly equicontinuous (see Exercise 1; Corollary 2 to Proposition 1; and no. 2, Proposition 4).

Examples. 1) Let X be a topological space (resp. a uniform space) and Y a uniform space. Every finite set of continuous (resp. uniformly continuous) mappings of X into Y is equicontinuous (resp. uniformly equicontinuous).

\begin{enumerate}
\def\labelenumi{\arabic{enumi})}
\setcounter{enumi}{1}
\tightlist
\item
  Let X, Y be two metric spaces, d (resp. d') the metric on X (resp. Y), and let k, \(\alpha\) be two real numbers \textgreater{} 0. Then the set of all mappings f: X → Y such that
\end{enumerate}

\[
d ^ {\prime} (f (x), f (x ^ {\prime})) \leqslant k (d (x, x ^ {\prime})) ^ {\alpha}
\]

for each pair of points \(x, x'\) of \(\mathbf{X}\) , is uniformly equicontinuous. For example, the set of all isometries (Chapter IX, § 2, no. 2) of \(\mathbf{X}\) onto a subset of \(\mathbf{Y}\) is uniformly equicontinuous.

* Let H be a set of real-valued functions defined on an interval I ⊂ R, which are differentiable on I and are such that \(|f'(x)| \leqslant k\) for all \(x \in I\) and all \(f \in H\) . Then H is uniformly equicontinuous, because if \(x_{1}, x_{2}\) are any two points of I we have \(|f(x_{1}) - f(x_{2})| \leqslant k |x_{1} - x_{2}|\) for each \(f \in H\) by the mean value theorem. *

\begin{enumerate}
\def\labelenumi{\arabic{enumi})}
\setcounter{enumi}{2}
\tightlist
\item
  Let G be a topological group, let Y be a uniform space and let \(f: \mathbf{G} \to \mathbf{Y}\) be a uniformly continuous mapping {[}G being endowed with its left uniformity (Chapter III, § 3, no. 1){]}. For each \(s \in \mathbf{G}\) , let \(f_s\) be the mapping \(x \to f(sx)\) of G into Y. Then the set of mappings \(f_s(s \in \mathbf{G})\) is uniformly equicontinuous, since the relation \(x^{-1}x' \in V\) is equivalent to \((sx)^{-1}(sx') \in V\) .
\end{enumerate}

PROPOSITION 1. Let T be a set, let S be a set of subsets of T, let Y be a uniform space, X a topological (resp. uniform) space, and let f be a mapping of \(\mathbf{T} \times \mathbf{X}\) into \(\mathbf{Y}\) . For each \(\mathbf{A} \in \mathfrak{S}\) , let \(\mathbf{H}_{\mathbf{A}} \subset \mathcal{F}(\mathbf{X}; \mathbf{Y})\) be the set of all mappings \(x \to f(t, x)\) as \(t\) runs through \(\mathbf{A}\) . Then the mapping \(x \to f(., x)\) of \(\mathbf{X}\) into \(\mathcal{F}_{\mathfrak{S}}(\mathbf{T}; \mathbf{Y})\) is continuous at a point \(x_0 \in \mathbf{X}\) (resp. uniformly continuous) if and only if the set \(\mathbf{H}_{\mathbf{A}}\) is equicontinuous at \(x_0\) (resp. uniformly equicontinuous) for all \(\mathbf{A} \in \mathfrak{S}\) .

Consider first the particular case where \(\mathfrak{S} = \{\mathrm{T}\}\) , i.e.~\(\mathcal{F}_{\mathfrak{S}}(\mathrm{T};\mathrm{Y}) = \mathcal{F}_u(\mathrm{T};\mathrm{Y})\) . For each entourage V of Y, the condition \((f(.,x),f(.,x'))\in W(\mathrm{V})\) signifies that \((f(t,x),f(t,x'))\in \mathrm{V}\) for all \(t\in \mathrm{T}\) . To say that \(x\to f(.,x)\) is continuous at \(x_0\) (resp. is uniformly continuous) is therefore equivalent to saying that, for each entourage V of Y, there is a neighbourhood U of \(x_0\) in X (resp. an entourage M of X) such that the relation \(x\in \mathrm{U}\) {[}resp. \((x,x')\in \mathrm{M}]\) implies \((f(t,x),f(t,x_0))\in \mathrm{V}\) {[}resp. \((f(t,x),f(t,x'))\in \mathrm{V}]\) for all \(t\in \mathrm{T}\) , and the proposition follows from Definitions 1 and 2. In the general case, we have to express that, for each \(\mathrm{A}\in \mathfrak{S}\) , the mapping \(x\rightarrow f(.,x)|\mathrm{A}\) of X into \(\mathcal{F}_u(\mathrm{A};\mathrm{Y})\) is continuous at \(x_0\) (resp. uniformly continuous), by virtue of § 1, no. 2; from what has been said, this is equivalent to saying that, for each \(\mathrm{A}\in \mathfrak{S}\) , \(\mathrm{H}_{\mathrm{A}}\) is equicontinuous at \(x_0\) (resp. uniformly equicontinuous).

Proposition 1 allows us to translate Definitions 1 and 2 into forms which are sometimes useful, by applying it to the case where \(\mathbf{T} = \mathbf{H}\) and \(f\) is the mapping \((h, x) \to h(x)\) of \(\mathbf{H} \times \mathbf{X}\) into \(\mathbf{Y}\) ; since \(f(., x)\) is the mapping \(h \to h(x)\) of \(\mathbf{H}\) into \(\mathbf{Y}\) , we see that:

COROLLARY 1. Let X be a topological (resp. uniform) space, Y a uniform space and H a subset of \(\mathcal{F}\) (X; Y). For each \(x \in \mathbf{X}\) , let \(\tilde{x}\) denote the mapping \(h \to h(x)\) of H into Y. Then H is equicontinuous at \(x_0\) (resp. uniformly equicontinuous) if and only if the mapping \(x \to \tilde{x}\) of X into the uniform space \(\mathcal{F}_u\) (H; Y) is continuous at \(x_0\) (resp. uniformly continuous).

In particular, if X is compact, every continuous mapping of X into \(\mathcal{F}_{u}(H;Y)\) is uniformly continuous (Chapter II, § 4, no. 1, Theorem 2). Therefore:

COROLLARY 2. Let X be a compact space, Y a uniform space. Then every equicontinuous subset of \(\mathcal{F}(X;Y)\) is uniformly equicontinuous.

Now suppose we have a set T, a topological space X, a uniform space Y and a mapping \(f: T \times X \to Y\) . Let \(\tilde{f}\) denote the mapping \(x \to f(., x)\) of X into \(\mathcal{F}_{u}(T; Y)\) , and let us consider the canonical mapping \(\theta: (t, g) \to g(t)\) of \(T \times \mathcal{F}_{u}(T; Y)\) into Y. It is clear that the diagram

\[
\mathrm{T} \times \mathrm{X} \xrightarrow {f} \mathrm{Y}
\]

\[
\iota_ {\mathrm{T}} \times \tilde {f} \searrow \nearrow \theta
\]

\[
\mathbf {T} \times \mathcal {F} _ {u} (\mathbf {T}; \mathbf {Y})
\]

(where \(\iota_{\mathbf{T}}\) is the identity mapping of T) is commutative. Suppose now that T is endowed with a topology and that, for each \(x \in \mathbf{X}\) , the mapping \(f(., x): t \to f(t, x)\) is continuous; we can then replace \(\mathcal{F}_{u}(\mathrm{T}; \mathrm{Y})\) by \(\mathcal{C}_{u}(\mathrm{T}; \mathrm{Y})\) in the above diagram. But we know that \(\theta\) is continuous from § 1, no. 6, Proposition 9; hence if \(\tilde{f}\) is continuous it follows that \(f\) is continuous. Since the continuity of \(\tilde{f}\) can be expressed with the help of Proposition 1, we obtain the following result:

COROLLARY 3. Let T, X be topological spaces, let Y be a uniform space and let f be a mapping of T × X into Y. Then f is continuous, provided that the following conditions are satisfied:

\begin{enumerate}
\def\labelenumi{\arabic{enumi})}
\item
  For each \(x \in \mathbf{X}\) , the partial mapping \(t \to f(t, x)\) is continuous.
\item
  As \(t\) runs through \(\mathbf{T}\) , the partial mappings \(x \to f(t, x)\) form an equicontinuous subset of \(\mathfrak{F}(\mathbf{X}; \mathbf{Y})\) .
\end{enumerate}

In particular, take T to be a subset H of \(\mathcal{F}(\mathbf{X};\mathbf{Y})\) and \(f\) to be the canonical mapping \((h,x)\to h(x)\) of \(\mathbf{H}\times \mathbf{X}\) into Y; condition 1) of Corollary 3 means that H is endowed with a topology finer than that of pointwise convergence, and condition 2) means that H is equicontinuous. Hence:

COROLLARY 4. Let X be a topological space, Y a uniform space, H an equicontinuous set of mappings of X into Y. If H is endowed with the topology of pointwise convergence, then the mapping \((h, x) \to h(x)\) of \(H \times X\) into Y is continuous.

More intuitively, this expresses the fact that if \(h \in \mathbf{H}\) converges pointwise to \(h_0 \in \mathbf{H}\) and if \(x \in \mathbf{X}\) converges to \(x_0\) , then \(h(x)\) converges to \(h_0(x_0)\) .

COROLLARY 5. Let X be a topological space, let Y, Z be two uniform spaces and let H be an equicontinuous set of mappings of Y into Z. If H, C(X; Y) and C(X; Z) are endowed with the topology of pointwise convergence, then the mapping \((u, v) \to u \circ v\) of H × C(X; Y) into C(X; Z) is continuous.

We have to show that, for each \(x \in \mathbf{X}\) , the mapping \((u, v) \to u(v(x))\) of \(\mathbf{H} \times \mathcal{C}(\mathbf{X}; \mathbf{Y})\) into \(\mathbf{Z}\) is continuous. Now \(v \to v(x)\) is continuous on \(\mathbf{H}\) ( \(\S 1\) , no. 2, Remark 6), and it follows from Corollary 4 that \((u, y) \to u(y)\) is a continuous mapping of \(\mathbf{H} \times \mathbf{Y}\) into \(\mathbf{Z}\) ; since \((u, v) \to u(v(x))\) is the composition of \((u, y) \to u(y)\) and \((u, v) \to (u, v(x))\) , the result is proved.

The following proposition and its corollary are the analogues of Corollaries 3 and 4 of Proposition 1 for uniformly equicontinuous sets of mappings:

PROPOSITION 2. Let T, X, Y be uniform spaces and let f be a mapping of T × X into Y. Then f is uniformly continuous if and only if the following two conditions are satisfied:

\begin{enumerate}
\def\labelenumi{\arabic{enumi})}
\item
  The mappings \(x \to f(t, x)\) ( \(t \in \mathbf{T}\) ) form a uniformly equicontinuous subset of \(\mathcal{F}(\mathbf{X}; \mathbf{Y})\) .
\item
  The mappings \(t \to f(t, x)\) ( \(x \in \mathbf{X}\) ) form a uniformly equicontinuous subset of \(\mathcal{F}(\mathbf{T}; \mathbf{Y})\) .
\end{enumerate}

It is easily seen that the conditions are necessary. Conversely, suppose that they are satisfied. Let W be an entourage of Y; then there exists an entourage U of T and an entourage V of X such that: 1) \((t', t'') \in \mathbf{U}\) implies that, for each \(x \in X\) ,

\[
\left(f \left(t ^ {\prime}, x\right), f \left(t ^ {\prime \prime}, x\right)\right) \in W.
\]

\begin{enumerate}
\def\labelenumi{\arabic{enumi})}
\setcounter{enumi}{1}
\tightlist
\item
  \((x', x'') \in \mathbf{V}\) implies that, for each \(t \in \mathbf{T}\) ,
\end{enumerate}

\[
(f (t, x ^ {\prime}), f (t, x ^ {\prime \prime})) \in W.
\]

It is now clear that the relation `` \((t', t'') \in \mathbf{U}\) and \((x', x'') \in \mathbf{V}\) '' implies that \((f(t', x'), f(t'', x'')) \in \overset{2}{\mathbf{W}}\) , whence the result.

In particular, take T to be a subset H of \(\mathcal{F}(\mathbf{X};\mathbf{Y})\) , endowed with the uniformity of uniform convergence, and take \(f\) to be the canonical mapping \((h,x)\to h(x)\) ; then condition 2) of Proposition 2 is automatically satisfied because, for each entourage W of Y, the set of pairs \((h',h'')\) such that \((h'(x),h''(x))\in W\) for all \(x\in X\) is by definition an entourage of the uniform structure of H. Hence only condition 1) has to be expressed; in other words:

COROLLARY. Let X, Y be two uniform spaces and let H be a subset of \(\mathcal{F}(X;Y)\) . Then H is uniformly equicontinuous if and only if the mapping \((h,x) \to h(x)\) of \(H \times X\) into Y is uniformly continuous, H being endowed with the uniformity of uniform convergence.

